\documentclass[aps,prl,reprint,superscriptaddress,nofootinbib,longbibliography]{revtex4-2}
\usepackage{amsmath,amssymb,amsthm,mathtools,bm,microtype}
\usepackage[hidelinks]{hyperref}
\newcommand{\Q}{\mathcal Q}\newcommand{\W}{\mathcal W}
\newcommand{\sxdeg}{\operatorname{sxdeg}}
\newtheorem{theorem}{Theorem}
\begin{document}
\title{Infinite NPA depth at minimal semidefinite complexity}
\author{Anubhav Chaturvedi}
\affiliation{Division of Quantum Optics and Information, Institute of Theoretical Physics and Astrophysics, Faculty of Mathematics, Physics and Informatics, University of Gda\'nsk, 80-308 Gda\'nsk, Poland}
\date{\today}
\begin{abstract}
The Navascu\'es--Pironio--Ac\'in hierarchy approaches quantum correlations through increasingly large semidefinite programs. We construct a quantum region that admits one fixed semidefinite lift although every finite hierarchy level fails to describe it. The region records the CHSH value and the sum of two marginal averages, so its supporting lines include the symmetric doubly tilted CHSH family. Equal measurement overlaps suffice for every supporting inequality: a determinant comparison handles the nonclassical directions, and the remaining bounds are local. The resulting univariate quartic matrix depends affinely on the inequality coefficients and yields an explicit lift with blocks of orders five and four. Independently varying the two measurement angles produces a shifted Motzkin polynomial that obstructs every fixed NPA level near the local threshold. The region has semidefinite extension degree two, the smallest possible for a nonpolyhedral set. Physical relabellings give thirty-two exact projections that can be added to any NPA relaxation.
\end{abstract}
\maketitle

Bell experiments ask which correlations can arise from local measurements
on a shared quantum state. The Navascu\'es--Pironio--Ac\'in (NPA)
hierarchy~\cite{NPA2007,NPA2008} answers through semidefinite outer
approximations: increasing the list of observable words gives
progressively stronger consistency conditions. The first level already
recovers the quantum CHSH bound~\cite{CHSH,Tsirelson}. Retaining marginal
averages changes the situation. Along the symmetric doubly tilted CHSH
family, every fixed level overestimates quantum bounds arbitrarily close
to the local threshold~\cite{prior}.

Does this unbounded depth reflect an intrinsically complicated quantum
region? We answer by constructing a fixed semidefinite representation
of the entire two-dimensional region probed by that family.
Its boundary includes all the directions responsible for NPA
nontermination. Moreover, a general planar theorem implies that this
region has the smallest possible semidefinite block order for a
nonpolyhedral set.

The construction turns on a reduction that preserves every support
value. A Bell operator initially depends on two independent measurement
overlaps. In the nonclassical range, its determinant is minimized at
fixed overlap product when the overlaps coincide. The remaining bounds
are attained by equal-overlap deterministic strategies. Thus every
supporting inequality can be decided with one variable. Positivity of a univariate
matrix polynomial then has a finite sum-of-squares representation.
The NPA obstruction retains both independent variables because its
operator certificates must hold before this support reduction.

\emph{The observed region.}---Alice and Bob each choose one of two
measurements with outcomes $\pm1$. Write
$a_x=\langle A_x\rangle$, $b_y=\langle B_y\rangle$ and
$E_{xy}=\langle A_xB_y\rangle$, and retain
\begin{equation}
 M=a_0+b_0,\qquad S=E_{00}+E_{01}+E_{10}-E_{11}.
 \label{eq:observed}
\end{equation}
Let $\Q_{222}$ be the quantum set in this scenario. The projection
\begin{equation}
 \W=\{(M,S):p\in\Q_{222}\}
 \label{eq:W}
\end{equation}
is compact and convex. It includes all quantum behaviors; equal
marginals are not assumed. Denote by $\W_L$ the same projection of the
standard level-$L$ NPA relaxation.

\begin{theorem}
The quantum region $\W$ has an exact affine semidefinite lift with
positive semidefinite blocks of orders five and four and sixteen
auxiliary variables. It satisfies
\begin{equation}
 \W\subsetneq\W_L\quad\text{for every fixed finite }L,
 \qquad \sxdeg(\W)=2 .
 \label{eq:theorem}
\end{equation}
\end{theorem}
Here $\sxdeg$ is the smallest possible largest block order among
finite semidefinite lifts. We first construct the explicit lift,
then connect it to the NPA obstruction. The optimal-order statement
follows at the end. Complete proofs and membership matrices are
provided in the Supplemental Material~\cite{SM}.

\emph{Equal overlaps determine all supporting lines.}---A pair
$(M,S)$ is determined by the inequalities
$hS+\alpha M\leq\lambda$ valid on $\W$, where $h,\alpha\in\mathbb R$.
Jordan's lemma reduces their validity to real two-qubit Bell operators
\begin{equation}
 H=h\,\mathrm{CHSH}+\alpha(Z\otimes I+I\otimes Z),
 \label{eq:H}
\end{equation}
with $A_0=B_0=Z$ and
\[
 A_1=cZ+\sqrt{1-c^2}X,\quad
 B_1=dZ+\sqrt{1-d^2}X.
\]
Here $c,d\in[-1,1]$; one-dimensional blocks are included at the endpoints.
Deterministic assignments attain
\begin{equation}
 \ell=2|h|+2|\alpha|,
 \label{eq:local}
\end{equation}
so a valid bound necessarily has $\lambda\geq\ell$.

On this range, the first three elementary symmetric functions of
the eigenvalues of $\lambda I-H$ are nonnegative~\cite{SM}.
The determinant therefore decides positivity. Indeed, if the
determinant is also nonnegative, the polynomial
$\det(rI+\lambda I-H)$ has nonnegative coefficients and cannot have
a positive root; a negative eigenvalue would supply such a root.

Set $x=(1-c)/2$, $y=(1-d)/2$ and
$D(x,y)=\det(\lambda I-H)$. Direct expansion gives
\begin{equation}
\begin{split}
 D(x,y)-D(\sqrt{xy},\sqrt{xy})
 ={}&64h^2(4h^2-\alpha^2)xy\\
 &{}\times(\sqrt{x}-\sqrt{y})^2 .
\end{split}
 \label{eq:det}
\end{equation}
For $|\alpha|\leq|h|$, the right-hand side is nonnegative.
Every pair of overlaps can thus be replaced by equal overlaps,
at the same $xy$, without increasing the determinant.
Checking all equal-overlap slacks proves the bound for every
measurement pair. This comparison is the reason the symmetric
optimization reduces to one angle.

The remaining directions have support $\ell$. To see this,
the one-sided tilted-CHSH inequality
$\mathrm{CHSH}+2A_0\leq4I$ follows from a sum of squares~\cite{SM}.
Party exchange and output flips give either marginal and either
sign of CHSH. Average the two appropriate inequalities, multiply
by $|h|$, and add the marginal slack of weight $|\alpha|-|h|$.
For $|\alpha|\geq|h|$ this proves $H\leq\ell I$.
Equal-overlap deterministic strategies attain $\ell$.
Thus the reduction covers all real support coefficients, including
zero coefficients and the critical line.

\emph{The finite lift.}---When $c=d$, the singlet decouples with
eigenvalue $-2hc$. The other three Bell states carry a real
$3\times3$ block. Parametrize
$c=(1-t^2)/(1+t^2)$ and multiply its slack by $(1+t^2)^2$.
Writing
\[
 p_0=(1+t^2)^2,\qquad p_1=1-t^4,\qquad p_2=1+4t^2-t^4,
\]
gives
\begin{equation}
 K(t)=
 \begin{pmatrix}
 \lambda p_0-2hp_1&-2\alpha p_0&0\\
 -2\alpha p_0&\lambda p_0-2hp_2&-8ht^3\\
 0&-8ht^3&\lambda p_0+2hp_2
 \end{pmatrix}.
 \label{eq:K}
\end{equation}
Positivity at $t=0$ and of the leading coefficient implies
$\lambda\geq\ell$, which also bounds the singlet.
The endpoint $c=-1$ follows by continuity. Negative $t$ adds no
condition since $K(-t)=RK(t)R$, with $R=\operatorname{diag}(1,1,-1)$.
We obtain
\begin{equation}
 \max_{(M,S)\in\W}(hS+\alpha M)\leq\lambda
 \quad\Longleftrightarrow\quad
 K(t)\succeq0\quad\forall t\in\mathbb R .
 \label{eq:equiv}
\end{equation}

Every PSD real univariate matrix polynomial is a polynomial matrix
sum of squares, including identically singular
polynomials~\cite{HanselkaSinn}. Since $K$ is quartic, its factors
have degree at most two: the highest-degree real squares cannot
cancel. Averaging a Gram representation under $t\mapsto-t$ and $R$
splits it into two parity blocks. For $\xi\in\mathbb R^3$, use
\begin{equation}
\begin{aligned}
 Z_+&=(\xi_1,\xi_2,t\xi_3,t^2\xi_1,t^2\xi_2)^T,\\
 Z_-&=(\xi_3,t\xi_1,t\xi_2,t^2\xi_3)^T .
\end{aligned}
\label{eq:features}
\end{equation}
Equation~\eqref{eq:equiv} is then equivalent to
\begin{equation}
 \xi^TK(t)\xi=Z_+^TGZ_++Z_-^TJZ_-,
 \qquad G\succeq0,\quad J\succeq0 .
 \label{eq:sos}
\end{equation}
The blocks have orders five and four. Coefficient matching gives
sixteen affine equations in $(G,J,\lambda,\alpha,h)$.

The affine dependence on \emph{all} support coefficients is essential:
these equations describe the complete cone of valid inequalities
of $\W$. Their dual supplies a single membership lift, with two
PSD matrices $U(z),V(z)$ and three affine equations
\begin{equation}
 \W=\left\{(M,S):
 \begin{array}{l}
 \exists z\in\mathbb R^{16},\ U(z),V(z)\succeq0,\\
 \ell(z)=1,\ m(z)=M,\ s(z)=S
 \end{array}\right\}.
 \label{eq:membership}
\end{equation}
All entries are rational and listed in the Supplemental Material.
Exact boundary membership follows from two elementary facts.
The Gram pair
\[
 G_0=\operatorname{diag}(1,1,2,1,1),\qquad
 J_0=\operatorname{diag}(1,2,2,1)
\]
represents $(\lambda,\alpha,h)=(1,0,0)$ and is positive definite.
It supplies strict feasibility, while
$\ell(z)=\operatorname{tr}(U(z)G_0)+\operatorname{tr}(V(z)J_0)$
bounds every dual variable on the normalized section.
The dual image is compact there. Conic duality and
Eq.~\eqref{eq:equiv} therefore prove both inclusions in
Eq.~\eqref{eq:membership}, with attained auxiliary variables.

\emph{The two-angle obstruction to finite NPA depth.}---For
$0\leq\alpha<1$, define the symmetric doubly tilted support
\[
 \omega_{\rm Q}(\alpha)=\max_{\W}(S+\alpha M),\qquad
 \Delta_\alpha=\omega_{\rm Q}(\alpha)-(2+2\alpha).
\]
The exact quantum bound and optimal measurements were obtained in
Ref.~\cite{GigenaEtAl2025}. Near the local threshold, $T=1-\alpha$
gives~\cite{prior,SM}
\begin{equation}
 \Delta_{1-T}=\frac43T^3+O(T^4).
 \label{eq:cubic}
\end{equation}
The NPA nontermination theorem states that, for every fixed finite $L$,
\begin{equation}
 \frac{\omega_{\rm NPA}^{(L)}(1-T)-\omega_{\rm Q}(1-T)}
 {\Delta_{1-T}}\longrightarrow+\infty
 \qquad(T\downarrow0).
 \label{eq:divergence}
\end{equation}
It concerns this very projection, so it implies the strict inclusion
in Eq.~\eqref{eq:theorem}.
Each fixed level's excess dominates the vanishing quantum advantage
above the local bound.

The proof identifies what remains difficult for NPA after the support
values have become univariate. Let $T=t^2$, keep $A_0=B_0=Z$, and
take the two other measurement angles independently as $tu,tv$.
Use the normalized version of
\begin{equation}
 \xi_t=(1,\ t^3u^2v/12,\ t^3uv^2/12,\ -t^2uv/4)^T .
 \label{eq:state}
\end{equation}
The gap between $\omega_{\rm Q}(1-t^2)$ and the value of this strategy,
divided by $\Delta_{1-t^2}$, converges locally uniformly to
\begin{equation}
 \mathcal M(u/2,v/2),\qquad
 \mathcal M(U,V)=1+U^4V^2+U^2V^4-3U^2V^2 .
 \label{eq:motzkin}
\end{equation}
This is the Motzkin polynomial.

Suppose the ratio in Eq.~\eqref{eq:divergence} stayed bounded along
a sequence. At a fixed NPA level, an optimal dual certificate is a
sum of squares in a fixed finite list of observable words.
A trace bound controls its Gram coefficients. Evaluate the square
factors on Eq.~\eqref{eq:state}, divide by $t^3$, and expand through
third order. The fixed word length gives a uniformly vanishing
Taylor remainder. Bounded sums of squared values then bound the
coefficients of these degree-three polynomials.
A convergent subsequence would give a polynomial sum of squares for
$\mathcal M+c$, with $c\geq0$.

Such squares cannot exist. Every factor has degree at most three;
restriction to either axis makes it constant on that axis. It
therefore has the form $a+bUV+dU^2V+eUV^2$.
The coefficient of $U^2V^2$ in the sum of squares is $\sum b^2\geq0$,
contradicting the coefficient $-3$ in $\mathcal M+c$.
The Supplemental Material supplies the trace, remainder and
coefficient bounds and handles arbitrary fixed finite word lists.

There is no conflict with Eq.~\eqref{eq:sos}. Equal overlaps determine
the quantum support values, whereas a finite NPA certificate must
represent the bound simultaneously for independently chosen
measurement angles. The polynomial squares obtained after
Eq.~\eqref{eq:det} need not be fixed-word operator identities.

\emph{Minimal block order and exact additions to NPA.}---Equation
\eqref{eq:membership} makes $\W$ a closed convex semialgebraic
subset of the plane. Scheiderer's theorem~\cite{Scheiderer} therefore
gives $\sxdeg(\W)\leq2$. It cannot be one, since order-one lifts
are polyhedral. Indeed, a polygon would have an affine support
function for $\alpha$ sufficiently close to one. Its limiting value
and slope would force the line $2+2\alpha$, contradicting the positive
cubic term in Eq.~\eqref{eq:cubic}. Thus $\sxdeg(\W)=2$.
This optimal-order conclusion is existential; Eqs.~\eqref{eq:sos}
and \eqref{eq:membership} provide the explicit blocks $5,4$.

The construction also gives exact information about many other
projections of the full quantum set. For each CHSH sign pattern
$\chi\in\{\pm1\}^{2\times2}$ with $\prod_{ij}\chi_{ij}=-1$, and each
selected pair $(x,y)$, define
\begin{equation}
 \pi_{\chi,xy}(p)=
 \left(a_x+\chi_{xy}b_y,\ \sum_{ij}\chi_{ij}E_{ij}\right).
 \label{eq:32}
\end{equation}
Physical relabellings show that each image
$\pi_{\chi,xy}(\Q_{222})$ is exactly $\W$.
The eight patterns and four selected pairs give 32 distinct planes.
Consequently
\begin{equation}
 \mathcal R_L=\Q_L\cap
 \bigcap_{\chi,x,y}\pi_{\chi,xy}^{-1}(\W)
 \label{eq:hybrid}
\end{equation}
retains every quantum behavior and has exact quantum support on
all these planes. It strictly strengthens each finite NPA relaxation.
The accompanying rational data implement the 32 added lifts without
any angle discretization.

\begin{thebibliography}{10}
\bibitem{NPA2007} M. Navascu\'es, S. Pironio, and A. Ac\'in,
Bounding the set of quantum correlations, Phys. Rev. Lett.
\textbf{98}, 010401 (2007).
\bibitem{NPA2008} M. Navascu\'es, S. Pironio, and A. Ac\'in,
A convergent hierarchy of semidefinite programs characterizing the set
of quantum correlations, New J. Phys. \textbf{10}, 073013 (2008).
\bibitem{CHSH} J. F. Clauser, M. A. Horne, A. Shimony, and R. A. Holt,
Proposed experiment to test local hidden-variable theories,
Phys. Rev. Lett. \textbf{23}, 880 (1969).
\bibitem{Tsirelson} B. S. Cirel'son, Quantum generalizations of Bell's
inequality, Lett. Math. Phys. \textbf{4}, 93 (1980).
\bibitem{prior} A. Chaturvedi, \emph{No Finite NPA Level Characterizes
the Complete Quantum Set in the Simplest Bell Scenario}, project
manuscript, September 2026. The complete nontermination argument
is reproduced in Part II of the Supplemental Material.
\bibitem{SM} See the accompanying Supplemental Material for the complete
finite-lift and NPA proofs, all signed asymmetric support branches,
and the physical relabelling construction.
\bibitem{HanselkaSinn} C. Hanselka and R. Sinn,
Positive semidefinite univariate matrix polynomials,
Math. Z. \textbf{292}, 83--101 (2019), Corollary 2.8.
\bibitem{GigenaEtAl2025} N. Gigena, E. Panwar, G. Scala, M. Ara\'ujo,
M. Farkas, and A. Chaturvedi, Self-testing tilted strategies for maximal
loophole-free nonlocality, npj Quantum Inf. \textbf{11}, 82 (2025).
\bibitem{Scheiderer} C. Scheiderer, Second-order cone representation
for convex sets in the plane, SIAM J. Appl. Algebra Geom.
\textbf{5}, 114--139 (2021), Theorem 0.1.
\end{thebibliography}
\end{document}
